\section{账本事件的区域史深描：liao}

\label{sec:ledger-depth-liao}

本组把账本事件置于战争、制度、区域社会与证据类型的交叉处。年表提供政治节点，地方材料与物质遗存则限制我们对征发、身份、信仰和迁徙所能作出的断言。

\subsection{战争、行政与区域资源}

账本条目\texttt{LIAO-0907-EVENT-026}所记907年的“耶律阿保机取得汗位，开始重组部族权力”，属于契丹的历史语境。条目所示前因是前期边防、财政和统治联盟的压力构成直接背景，具体条件须按本条材料分辨，其可见后果为改变后续资源配置与政治选择；实际执行范围和社会影响仍须分项核验。这一变化若进入区域生活，必须通过道路、文书、军队、市场、家庭和地方中介才会发生；政权易手或法令发布不等于所有城镇、乡村与边地在同一时间获得同样经验。

本条所列来源为《辽史》〈本纪〉、〈营卫志〉、〈百官志〉；《宋史》辽国相关传；宋人使辽记录。它们能够支持特定的年序、制度语言、政治行动或局部遗存，却不能无条件化为社会总量。账本特别提醒“编年主干可靠；细节、规模与动机须与对方材料、地方文书或考古材料互校”。因此，军报的兵数、条约的名分、官府的族群分类、判牍中的道德判断以及后出的叙述都须回到产生它们的场景，而不可直接代表全域动机或人口。

将事件与地方层次连接时，应同时考察资源怎样被征集、信息怎样传递、迁徙怎样改变户籍和劳动、宗教与亲属网络怎样提供或限制协作。现有证据往往只照亮少数地点；保留这种不均衡，才能避免把宋辽夏金元的多政权历史改写成单一中心的必然过程。

账本条目\texttt{LIAO-0907-REGNAL-YEAR}所记907年的“契丹／辽以辽太祖在位第1年作为诏令、赋役与外交文书的纪年框架”，属于契丹／辽的历史语境。条目所示前因是新岁颁历、年号和纪年是中枢与地方行政沟通的共同前提，其可见后果为为同年政令、战事和地方材料提供日期锚点，不能单独推知具体政策执行。这一变化若进入区域生活，必须通过道路、文书、军队、市场、家庭和地方中介才会发生；政权易手或法令发布不等于所有城镇、乡村与边地在同一时间获得同样经验。

本条所列来源为《辽史》〈本纪〉、〈营卫志〉、〈百官志〉；《宋史》辽国相关传；宋人使辽记录。它们能够支持特定的年序、制度语言、政治行动或局部遗存，却不能无条件化为社会总量。账本特别提醒“纪年框架可靠；该记录仅确证政权纪年连续，年内具体事项须由本年条另证”。因此，军报的兵数、条约的名分、官府的族群分类、判牍中的道德判断以及后出的叙述都须回到产生它们的场景，而不可直接代表全域动机或人口。

将事件与地方层次连接时，应同时考察资源怎样被征集、信息怎样传递、迁徙怎样改变户籍和劳动、宗教与亲属网络怎样提供或限制协作。现有证据往往只照亮少数地点；保留这种不均衡，才能避免把宋辽夏金元的多政权历史改写成单一中心的必然过程。

账本条目\texttt{LIAO-0908-REGNAL-YEAR}所记908年的“契丹／辽以辽太祖在位第2年作为诏令、赋役与外交文书的纪年框架”，属于契丹／辽的历史语境。条目所示前因是新岁颁历、年号和纪年是中枢与地方行政沟通的共同前提，其可见后果为为同年政令、战事和地方材料提供日期锚点，不能单独推知具体政策执行。这一变化若进入区域生活，必须通过道路、文书、军队、市场、家庭和地方中介才会发生；政权易手或法令发布不等于所有城镇、乡村与边地在同一时间获得同样经验。

本条所列来源为《辽史》〈本纪〉、〈营卫志〉、〈百官志〉；《宋史》辽国相关传；宋人使辽记录。它们能够支持特定的年序、制度语言、政治行动或局部遗存，却不能无条件化为社会总量。账本特别提醒“纪年框架可靠；该记录仅确证政权纪年连续，年内具体事项须由本年条另证”。因此，军报的兵数、条约的名分、官府的族群分类、判牍中的道德判断以及后出的叙述都须回到产生它们的场景，而不可直接代表全域动机或人口。

将事件与地方层次连接时，应同时考察资源怎样被征集、信息怎样传递、迁徙怎样改变户籍和劳动、宗教与亲属网络怎样提供或限制协作。现有证据往往只照亮少数地点；保留这种不均衡，才能避免把宋辽夏金元的多政权历史改写成单一中心的必然过程。

账本条目\texttt{LIAO-0909-REGNAL-YEAR}所记909年的“契丹／辽以辽太祖在位第3年作为诏令、赋役与外交文书的纪年框架”，属于契丹／辽的历史语境。条目所示前因是新岁颁历、年号和纪年是中枢与地方行政沟通的共同前提，其可见后果为为同年政令、战事和地方材料提供日期锚点，不能单独推知具体政策执行。这一变化若进入区域生活，必须通过道路、文书、军队、市场、家庭和地方中介才会发生；政权易手或法令发布不等于所有城镇、乡村与边地在同一时间获得同样经验。

本条所列来源为《辽史》〈本纪〉、〈营卫志〉、〈百官志〉；《宋史》辽国相关传；宋人使辽记录。它们能够支持特定的年序、制度语言、政治行动或局部遗存，却不能无条件化为社会总量。账本特别提醒“纪年框架可靠；该记录仅确证政权纪年连续，年内具体事项须由本年条另证”。因此，军报的兵数、条约的名分、官府的族群分类、判牍中的道德判断以及后出的叙述都须回到产生它们的场景，而不可直接代表全域动机或人口。

将事件与地方层次连接时，应同时考察资源怎样被征集、信息怎样传递、迁徙怎样改变户籍和劳动、宗教与亲属网络怎样提供或限制协作。现有证据往往只照亮少数地点；保留这种不均衡，才能避免把宋辽夏金元的多政权历史改写成单一中心的必然过程。

账本条目\texttt{LIAO-0910-REGNAL-YEAR}所记910年的“契丹／辽以辽太祖在位第4年作为诏令、赋役与外交文书的纪年框架”，属于契丹／辽的历史语境。条目所示前因是新岁颁历、年号和纪年是中枢与地方行政沟通的共同前提，其可见后果为为同年政令、战事和地方材料提供日期锚点，不能单独推知具体政策执行。这一变化若进入区域生活，必须通过道路、文书、军队、市场、家庭和地方中介才会发生；政权易手或法令发布不等于所有城镇、乡村与边地在同一时间获得同样经验。

本条所列来源为《辽史》〈本纪〉、〈营卫志〉、〈百官志〉；《宋史》辽国相关传；宋人使辽记录。它们能够支持特定的年序、制度语言、政治行动或局部遗存，却不能无条件化为社会总量。账本特别提醒“纪年框架可靠；该记录仅确证政权纪年连续，年内具体事项须由本年条另证”。因此，军报的兵数、条约的名分、官府的族群分类、判牍中的道德判断以及后出的叙述都须回到产生它们的场景，而不可直接代表全域动机或人口。

将事件与地方层次连接时，应同时考察资源怎样被征集、信息怎样传递、迁徙怎样改变户籍和劳动、宗教与亲属网络怎样提供或限制协作。现有证据往往只照亮少数地点；保留这种不均衡，才能避免把宋辽夏金元的多政权历史改写成单一中心的必然过程。

账本条目\texttt{LIAO-0911-REGNAL-YEAR}所记911年的“契丹／辽以辽太祖在位第5年作为诏令、赋役与外交文书的纪年框架”，属于契丹／辽的历史语境。条目所示前因是新岁颁历、年号和纪年是中枢与地方行政沟通的共同前提，其可见后果为为同年政令、战事和地方材料提供日期锚点，不能单独推知具体政策执行。这一变化若进入区域生活，必须通过道路、文书、军队、市场、家庭和地方中介才会发生；政权易手或法令发布不等于所有城镇、乡村与边地在同一时间获得同样经验。

本条所列来源为《辽史》〈本纪〉、〈营卫志〉、〈百官志〉；《宋史》辽国相关传；宋人使辽记录。它们能够支持特定的年序、制度语言、政治行动或局部遗存，却不能无条件化为社会总量。账本特别提醒“纪年框架可靠；该记录仅确证政权纪年连续，年内具体事项须由本年条另证”。因此，军报的兵数、条约的名分、官府的族群分类、判牍中的道德判断以及后出的叙述都须回到产生它们的场景，而不可直接代表全域动机或人口。

将事件与地方层次连接时，应同时考察资源怎样被征集、信息怎样传递、迁徙怎样改变户籍和劳动、宗教与亲属网络怎样提供或限制协作。现有证据往往只照亮少数地点；保留这种不均衡，才能避免把宋辽夏金元的多政权历史改写成单一中心的必然过程。

账本条目\texttt{LIAO-0912-REGNAL-YEAR}所记912年的“契丹／辽以辽太祖在位第6年作为诏令、赋役与外交文书的纪年框架”，属于契丹／辽的历史语境。条目所示前因是新岁颁历、年号和纪年是中枢与地方行政沟通的共同前提，其可见后果为为同年政令、战事和地方材料提供日期锚点，不能单独推知具体政策执行。这一变化若进入区域生活，必须通过道路、文书、军队、市场、家庭和地方中介才会发生；政权易手或法令发布不等于所有城镇、乡村与边地在同一时间获得同样经验。

本条所列来源为《辽史》〈本纪〉、〈营卫志〉、〈百官志〉；《宋史》辽国相关传；宋人使辽记录。它们能够支持特定的年序、制度语言、政治行动或局部遗存，却不能无条件化为社会总量。账本特别提醒“纪年框架可靠；该记录仅确证政权纪年连续，年内具体事项须由本年条另证”。因此，军报的兵数、条约的名分、官府的族群分类、判牍中的道德判断以及后出的叙述都须回到产生它们的场景，而不可直接代表全域动机或人口。

将事件与地方层次连接时，应同时考察资源怎样被征集、信息怎样传递、迁徙怎样改变户籍和劳动、宗教与亲属网络怎样提供或限制协作。现有证据往往只照亮少数地点；保留这种不均衡，才能避免把宋辽夏金元的多政权历史改写成单一中心的必然过程。

账本条目\texttt{LIAO-0913-REGNAL-YEAR}所记913年的“契丹／辽以辽太祖在位第7年作为诏令、赋役与外交文书的纪年框架”，属于契丹／辽的历史语境。条目所示前因是新岁颁历、年号和纪年是中枢与地方行政沟通的共同前提，其可见后果为为同年政令、战事和地方材料提供日期锚点，不能单独推知具体政策执行。这一变化若进入区域生活，必须通过道路、文书、军队、市场、家庭和地方中介才会发生；政权易手或法令发布不等于所有城镇、乡村与边地在同一时间获得同样经验。

本条所列来源为《辽史》〈本纪〉、〈营卫志〉、〈百官志〉；《宋史》辽国相关传；宋人使辽记录。它们能够支持特定的年序、制度语言、政治行动或局部遗存，却不能无条件化为社会总量。账本特别提醒“纪年框架可靠；该记录仅确证政权纪年连续，年内具体事项须由本年条另证”。因此，军报的兵数、条约的名分、官府的族群分类、判牍中的道德判断以及后出的叙述都须回到产生它们的场景，而不可直接代表全域动机或人口。

将事件与地方层次连接时，应同时考察资源怎样被征集、信息怎样传递、迁徙怎样改变户籍和劳动、宗教与亲属网络怎样提供或限制协作。现有证据往往只照亮少数地点；保留这种不均衡，才能避免把宋辽夏金元的多政权历史改写成单一中心的必然过程。

账本条目\texttt{LIAO-0914-REGNAL-YEAR}所记914年的“契丹／辽以辽太祖在位第8年作为诏令、赋役与外交文书的纪年框架”，属于契丹／辽的历史语境。条目所示前因是新岁颁历、年号和纪年是中枢与地方行政沟通的共同前提，其可见后果为为同年政令、战事和地方材料提供日期锚点，不能单独推知具体政策执行。这一变化若进入区域生活，必须通过道路、文书、军队、市场、家庭和地方中介才会发生；政权易手或法令发布不等于所有城镇、乡村与边地在同一时间获得同样经验。

本条所列来源为《辽史》〈本纪〉、〈营卫志〉、〈百官志〉；《宋史》辽国相关传；宋人使辽记录。它们能够支持特定的年序、制度语言、政治行动或局部遗存，却不能无条件化为社会总量。账本特别提醒“纪年框架可靠；该记录仅确证政权纪年连续，年内具体事项须由本年条另证”。因此，军报的兵数、条约的名分、官府的族群分类、判牍中的道德判断以及后出的叙述都须回到产生它们的场景，而不可直接代表全域动机或人口。

将事件与地方层次连接时，应同时考察资源怎样被征集、信息怎样传递、迁徙怎样改变户籍和劳动、宗教与亲属网络怎样提供或限制协作。现有证据往往只照亮少数地点；保留这种不均衡，才能避免把宋辽夏金元的多政权历史改写成单一中心的必然过程。

账本条目\texttt{LIAO-0915-REGNAL-YEAR}所记915年的“契丹／辽以辽太祖在位第9年作为诏令、赋役与外交文书的纪年框架”，属于契丹／辽的历史语境。条目所示前因是新岁颁历、年号和纪年是中枢与地方行政沟通的共同前提，其可见后果为为同年政令、战事和地方材料提供日期锚点，不能单独推知具体政策执行。这一变化若进入区域生活，必须通过道路、文书、军队、市场、家庭和地方中介才会发生；政权易手或法令发布不等于所有城镇、乡村与边地在同一时间获得同样经验。

本条所列来源为《辽史》〈本纪〉、〈营卫志〉、〈百官志〉；《宋史》辽国相关传；宋人使辽记录。它们能够支持特定的年序、制度语言、政治行动或局部遗存，却不能无条件化为社会总量。账本特别提醒“纪年框架可靠；该记录仅确证政权纪年连续，年内具体事项须由本年条另证”。因此，军报的兵数、条约的名分、官府的族群分类、判牍中的道德判断以及后出的叙述都须回到产生它们的场景，而不可直接代表全域动机或人口。

将事件与地方层次连接时，应同时考察资源怎样被征集、信息怎样传递、迁徙怎样改变户籍和劳动、宗教与亲属网络怎样提供或限制协作。现有证据往往只照亮少数地点；保留这种不均衡，才能避免把宋辽夏金元的多政权历史改写成单一中心的必然过程。

账本条目\texttt{LIAO-0916-EVENT-027}所记916年的“阿保机称帝并建国号契丹”，属于契丹的历史语境。条目所示前因是前期边防、财政和统治联盟的压力构成直接背景，具体条件须按本条材料分辨，其可见后果为改变后续资源配置与政治选择；实际执行范围和社会影响仍须分项核验。这一变化若进入区域生活，必须通过道路、文书、军队、市场、家庭和地方中介才会发生；政权易手或法令发布不等于所有城镇、乡村与边地在同一时间获得同样经验。

本条所列来源为《辽史》〈本纪〉、〈营卫志〉、〈百官志〉；《宋史》辽国相关传；宋人使辽记录。它们能够支持特定的年序、制度语言、政治行动或局部遗存，却不能无条件化为社会总量。账本特别提醒“编年主干可靠；细节、规模与动机须与对方材料、地方文书或考古材料互校”。因此，军报的兵数、条约的名分、官府的族群分类、判牍中的道德判断以及后出的叙述都须回到产生它们的场景，而不可直接代表全域动机或人口。

将事件与地方层次连接时，应同时考察资源怎样被征集、信息怎样传递、迁徙怎样改变户籍和劳动、宗教与亲属网络怎样提供或限制协作。现有证据往往只照亮少数地点；保留这种不均衡，才能避免把宋辽夏金元的多政权历史改写成单一中心的必然过程。

账本条目\texttt{LIAO-0916-REGNAL-YEAR}所记916年的“契丹／辽以辽太祖在位第10年作为诏令、赋役与外交文书的纪年框架”，属于契丹／辽的历史语境。条目所示前因是新岁颁历、年号和纪年是中枢与地方行政沟通的共同前提，其可见后果为为同年政令、战事和地方材料提供日期锚点，不能单独推知具体政策执行。这一变化若进入区域生活，必须通过道路、文书、军队、市场、家庭和地方中介才会发生；政权易手或法令发布不等于所有城镇、乡村与边地在同一时间获得同样经验。

本条所列来源为《辽史》〈本纪〉、〈营卫志〉、〈百官志〉；《宋史》辽国相关传；宋人使辽记录。它们能够支持特定的年序、制度语言、政治行动或局部遗存，却不能无条件化为社会总量。账本特别提醒“纪年框架可靠；该记录仅确证政权纪年连续，年内具体事项须由本年条另证”。因此，军报的兵数、条约的名分、官府的族群分类、判牍中的道德判断以及后出的叙述都须回到产生它们的场景，而不可直接代表全域动机或人口。

将事件与地方层次连接时，应同时考察资源怎样被征集、信息怎样传递、迁徙怎样改变户籍和劳动、宗教与亲属网络怎样提供或限制协作。现有证据往往只照亮少数地点；保留这种不均衡，才能避免把宋辽夏金元的多政权历史改写成单一中心的必然过程。

\subsection{外交、道路与人口移动}

账本条目\texttt{LIAO-0917-REGNAL-YEAR}所记917年的“契丹／辽以辽太祖在位第11年作为诏令、赋役与外交文书的纪年框架”，属于契丹／辽的历史语境。条目所示前因是新岁颁历、年号和纪年是中枢与地方行政沟通的共同前提，其可见后果为为同年政令、战事和地方材料提供日期锚点，不能单独推知具体政策执行。这一变化若进入区域生活，必须通过道路、文书、军队、市场、家庭和地方中介才会发生；政权易手或法令发布不等于所有城镇、乡村与边地在同一时间获得同样经验。

本条所列来源为《辽史》〈本纪〉、〈营卫志〉、〈百官志〉；《宋史》辽国相关传；宋人使辽记录。它们能够支持特定的年序、制度语言、政治行动或局部遗存，却不能无条件化为社会总量。账本特别提醒“纪年框架可靠；该记录仅确证政权纪年连续，年内具体事项须由本年条另证”。因此，军报的兵数、条约的名分、官府的族群分类、判牍中的道德判断以及后出的叙述都须回到产生它们的场景，而不可直接代表全域动机或人口。

将事件与地方层次连接时，应同时考察资源怎样被征集、信息怎样传递、迁徙怎样改变户籍和劳动、宗教与亲属网络怎样提供或限制协作。现有证据往往只照亮少数地点；保留这种不均衡，才能避免把宋辽夏金元的多政权历史改写成单一中心的必然过程。

账本条目\texttt{LIAO-0918-REGNAL-YEAR}所记918年的“契丹／辽以辽太祖在位第12年作为诏令、赋役与外交文书的纪年框架”，属于契丹／辽的历史语境。条目所示前因是新岁颁历、年号和纪年是中枢与地方行政沟通的共同前提，其可见后果为为同年政令、战事和地方材料提供日期锚点，不能单独推知具体政策执行。这一变化若进入区域生活，必须通过道路、文书、军队、市场、家庭和地方中介才会发生；政权易手或法令发布不等于所有城镇、乡村与边地在同一时间获得同样经验。

本条所列来源为《辽史》〈本纪〉、〈营卫志〉、〈百官志〉；《宋史》辽国相关传；宋人使辽记录。它们能够支持特定的年序、制度语言、政治行动或局部遗存，却不能无条件化为社会总量。账本特别提醒“纪年框架可靠；该记录仅确证政权纪年连续，年内具体事项须由本年条另证”。因此，军报的兵数、条约的名分、官府的族群分类、判牍中的道德判断以及后出的叙述都须回到产生它们的场景，而不可直接代表全域动机或人口。

将事件与地方层次连接时，应同时考察资源怎样被征集、信息怎样传递、迁徙怎样改变户籍和劳动、宗教与亲属网络怎样提供或限制协作。现有证据往往只照亮少数地点；保留这种不均衡，才能避免把宋辽夏金元的多政权历史改写成单一中心的必然过程。

账本条目\texttt{LIAO-0919-REGNAL-YEAR}所记919年的“契丹／辽以辽太祖在位第13年作为诏令、赋役与外交文书的纪年框架”，属于契丹／辽的历史语境。条目所示前因是新岁颁历、年号和纪年是中枢与地方行政沟通的共同前提，其可见后果为为同年政令、战事和地方材料提供日期锚点，不能单独推知具体政策执行。这一变化若进入区域生活，必须通过道路、文书、军队、市场、家庭和地方中介才会发生；政权易手或法令发布不等于所有城镇、乡村与边地在同一时间获得同样经验。

本条所列来源为《辽史》〈本纪〉、〈营卫志〉、〈百官志〉；《宋史》辽国相关传；宋人使辽记录。它们能够支持特定的年序、制度语言、政治行动或局部遗存，却不能无条件化为社会总量。账本特别提醒“纪年框架可靠；该记录仅确证政权纪年连续，年内具体事项须由本年条另证”。因此，军报的兵数、条约的名分、官府的族群分类、判牍中的道德判断以及后出的叙述都须回到产生它们的场景，而不可直接代表全域动机或人口。

将事件与地方层次连接时，应同时考察资源怎样被征集、信息怎样传递、迁徙怎样改变户籍和劳动、宗教与亲属网络怎样提供或限制协作。现有证据往往只照亮少数地点；保留这种不均衡，才能避免把宋辽夏金元的多政权历史改写成单一中心的必然过程。

账本条目\texttt{LIAO-0920-REGNAL-YEAR}所记920年的“契丹／辽以辽太祖在位第14年作为诏令、赋役与外交文书的纪年框架”，属于契丹／辽的历史语境。条目所示前因是新岁颁历、年号和纪年是中枢与地方行政沟通的共同前提，其可见后果为为同年政令、战事和地方材料提供日期锚点，不能单独推知具体政策执行。这一变化若进入区域生活，必须通过道路、文书、军队、市场、家庭和地方中介才会发生；政权易手或法令发布不等于所有城镇、乡村与边地在同一时间获得同样经验。

本条所列来源为《辽史》〈本纪〉、〈营卫志〉、〈百官志〉；《宋史》辽国相关传；宋人使辽记录。它们能够支持特定的年序、制度语言、政治行动或局部遗存，却不能无条件化为社会总量。账本特别提醒“纪年框架可靠；该记录仅确证政权纪年连续，年内具体事项须由本年条另证”。因此，军报的兵数、条约的名分、官府的族群分类、判牍中的道德判断以及后出的叙述都须回到产生它们的场景，而不可直接代表全域动机或人口。

将事件与地方层次连接时，应同时考察资源怎样被征集、信息怎样传递、迁徙怎样改变户籍和劳动、宗教与亲属网络怎样提供或限制协作。现有证据往往只照亮少数地点；保留这种不均衡，才能避免把宋辽夏金元的多政权历史改写成单一中心的必然过程。

账本条目\texttt{LIAO-0921-REGNAL-YEAR}所记921年的“契丹／辽以辽太祖在位第15年作为诏令、赋役与外交文书的纪年框架”，属于契丹／辽的历史语境。条目所示前因是新岁颁历、年号和纪年是中枢与地方行政沟通的共同前提，其可见后果为为同年政令、战事和地方材料提供日期锚点，不能单独推知具体政策执行。这一变化若进入区域生活，必须通过道路、文书、军队、市场、家庭和地方中介才会发生；政权易手或法令发布不等于所有城镇、乡村与边地在同一时间获得同样经验。

本条所列来源为《辽史》〈本纪〉、〈营卫志〉、〈百官志〉；《宋史》辽国相关传；宋人使辽记录。它们能够支持特定的年序、制度语言、政治行动或局部遗存，却不能无条件化为社会总量。账本特别提醒“纪年框架可靠；该记录仅确证政权纪年连续，年内具体事项须由本年条另证”。因此，军报的兵数、条约的名分、官府的族群分类、判牍中的道德判断以及后出的叙述都须回到产生它们的场景，而不可直接代表全域动机或人口。

将事件与地方层次连接时，应同时考察资源怎样被征集、信息怎样传递、迁徙怎样改变户籍和劳动、宗教与亲属网络怎样提供或限制协作。现有证据往往只照亮少数地点；保留这种不均衡，才能避免把宋辽夏金元的多政权历史改写成单一中心的必然过程。

账本条目\texttt{LIAO-0922-REGNAL-YEAR}所记922年的“契丹／辽以辽太祖在位第16年作为诏令、赋役与外交文书的纪年框架”，属于契丹／辽的历史语境。条目所示前因是新岁颁历、年号和纪年是中枢与地方行政沟通的共同前提，其可见后果为为同年政令、战事和地方材料提供日期锚点，不能单独推知具体政策执行。这一变化若进入区域生活，必须通过道路、文书、军队、市场、家庭和地方中介才会发生；政权易手或法令发布不等于所有城镇、乡村与边地在同一时间获得同样经验。

本条所列来源为《辽史》〈本纪〉、〈营卫志〉、〈百官志〉；《宋史》辽国相关传；宋人使辽记录。它们能够支持特定的年序、制度语言、政治行动或局部遗存，却不能无条件化为社会总量。账本特别提醒“纪年框架可靠；该记录仅确证政权纪年连续，年内具体事项须由本年条另证”。因此，军报的兵数、条约的名分、官府的族群分类、判牍中的道德判断以及后出的叙述都须回到产生它们的场景，而不可直接代表全域动机或人口。

将事件与地方层次连接时，应同时考察资源怎样被征集、信息怎样传递、迁徙怎样改变户籍和劳动、宗教与亲属网络怎样提供或限制协作。现有证据往往只照亮少数地点；保留这种不均衡，才能避免把宋辽夏金元的多政权历史改写成单一中心的必然过程。

账本条目\texttt{LIAO-0923-REGNAL-YEAR}所记923年的“契丹／辽以辽太祖在位第17年作为诏令、赋役与外交文书的纪年框架”，属于契丹／辽的历史语境。条目所示前因是新岁颁历、年号和纪年是中枢与地方行政沟通的共同前提，其可见后果为为同年政令、战事和地方材料提供日期锚点，不能单独推知具体政策执行。这一变化若进入区域生活，必须通过道路、文书、军队、市场、家庭和地方中介才会发生；政权易手或法令发布不等于所有城镇、乡村与边地在同一时间获得同样经验。

本条所列来源为《辽史》〈本纪〉、〈营卫志〉、〈百官志〉；《宋史》辽国相关传；宋人使辽记录。它们能够支持特定的年序、制度语言、政治行动或局部遗存，却不能无条件化为社会总量。账本特别提醒“纪年框架可靠；该记录仅确证政权纪年连续，年内具体事项须由本年条另证”。因此，军报的兵数、条约的名分、官府的族群分类、判牍中的道德判断以及后出的叙述都须回到产生它们的场景，而不可直接代表全域动机或人口。

将事件与地方层次连接时，应同时考察资源怎样被征集、信息怎样传递、迁徙怎样改变户籍和劳动、宗教与亲属网络怎样提供或限制协作。现有证据往往只照亮少数地点；保留这种不均衡，才能避免把宋辽夏金元的多政权历史改写成单一中心的必然过程。

账本条目\texttt{LIAO-0924-REGNAL-YEAR}所记924年的“契丹／辽以辽太祖在位第18年作为诏令、赋役与外交文书的纪年框架”，属于契丹／辽的历史语境。条目所示前因是新岁颁历、年号和纪年是中枢与地方行政沟通的共同前提，其可见后果为为同年政令、战事和地方材料提供日期锚点，不能单独推知具体政策执行。这一变化若进入区域生活，必须通过道路、文书、军队、市场、家庭和地方中介才会发生；政权易手或法令发布不等于所有城镇、乡村与边地在同一时间获得同样经验。

本条所列来源为《辽史》〈本纪〉、〈营卫志〉、〈百官志〉；《宋史》辽国相关传；宋人使辽记录。它们能够支持特定的年序、制度语言、政治行动或局部遗存，却不能无条件化为社会总量。账本特别提醒“纪年框架可靠；该记录仅确证政权纪年连续，年内具体事项须由本年条另证”。因此，军报的兵数、条约的名分、官府的族群分类、判牍中的道德判断以及后出的叙述都须回到产生它们的场景，而不可直接代表全域动机或人口。

将事件与地方层次连接时，应同时考察资源怎样被征集、信息怎样传递、迁徙怎样改变户籍和劳动、宗教与亲属网络怎样提供或限制协作。现有证据往往只照亮少数地点；保留这种不均衡，才能避免把宋辽夏金元的多政权历史改写成单一中心的必然过程。

账本条目\texttt{LIAO-0925-REGNAL-YEAR}所记925年的“契丹／辽以辽太祖在位第19年作为诏令、赋役与外交文书的纪年框架”，属于契丹／辽的历史语境。条目所示前因是新岁颁历、年号和纪年是中枢与地方行政沟通的共同前提，其可见后果为为同年政令、战事和地方材料提供日期锚点，不能单独推知具体政策执行。这一变化若进入区域生活，必须通过道路、文书、军队、市场、家庭和地方中介才会发生；政权易手或法令发布不等于所有城镇、乡村与边地在同一时间获得同样经验。

本条所列来源为《辽史》〈本纪〉、〈营卫志〉、〈百官志〉；《宋史》辽国相关传；宋人使辽记录。它们能够支持特定的年序、制度语言、政治行动或局部遗存，却不能无条件化为社会总量。账本特别提醒“纪年框架可靠；该记录仅确证政权纪年连续，年内具体事项须由本年条另证”。因此，军报的兵数、条约的名分、官府的族群分类、判牍中的道德判断以及后出的叙述都须回到产生它们的场景，而不可直接代表全域动机或人口。

将事件与地方层次连接时，应同时考察资源怎样被征集、信息怎样传递、迁徙怎样改变户籍和劳动、宗教与亲属网络怎样提供或限制协作。现有证据往往只照亮少数地点；保留这种不均衡，才能避免把宋辽夏金元的多政权历史改写成单一中心的必然过程。

账本条目\texttt{LIAO-0926-EVENT-028}所记926年的“辽军攻灭渤海，设置东丹等安排”，属于辽与渤海的历史语境。条目所示前因是前期边防、财政和统治联盟的压力构成直接背景，具体条件须按本条材料分辨，其可见后果为改变后续资源配置与政治选择；实际执行范围和社会影响仍须分项核验。这一变化若进入区域生活，必须通过道路、文书、军队、市场、家庭和地方中介才会发生；政权易手或法令发布不等于所有城镇、乡村与边地在同一时间获得同样经验。

本条所列来源为《辽史》〈本纪〉、〈营卫志〉、〈百官志〉；《宋史》辽国相关传；宋人使辽记录。它们能够支持特定的年序、制度语言、政治行动或局部遗存，却不能无条件化为社会总量。账本特别提醒“编年主干可靠；细节、规模与动机须与对方材料、地方文书或考古材料互校”。因此，军报的兵数、条约的名分、官府的族群分类、判牍中的道德判断以及后出的叙述都须回到产生它们的场景，而不可直接代表全域动机或人口。

将事件与地方层次连接时，应同时考察资源怎样被征集、信息怎样传递、迁徙怎样改变户籍和劳动、宗教与亲属网络怎样提供或限制协作。现有证据往往只照亮少数地点；保留这种不均衡，才能避免把宋辽夏金元的多政权历史改写成单一中心的必然过程。

账本条目\texttt{LIAO-0926-REGNAL-YEAR}所记926年的“契丹／辽以辽太祖在位第20年作为诏令、赋役与外交文书的纪年框架”，属于契丹／辽的历史语境。条目所示前因是新岁颁历、年号和纪年是中枢与地方行政沟通的共同前提，其可见后果为为同年政令、战事和地方材料提供日期锚点，不能单独推知具体政策执行。这一变化若进入区域生活，必须通过道路、文书、军队、市场、家庭和地方中介才会发生；政权易手或法令发布不等于所有城镇、乡村与边地在同一时间获得同样经验。

本条所列来源为《辽史》〈本纪〉、〈营卫志〉、〈百官志〉；《宋史》辽国相关传；宋人使辽记录。它们能够支持特定的年序、制度语言、政治行动或局部遗存，却不能无条件化为社会总量。账本特别提醒“纪年框架可靠；该记录仅确证政权纪年连续，年内具体事项须由本年条另证”。因此，军报的兵数、条约的名分、官府的族群分类、判牍中的道德判断以及后出的叙述都须回到产生它们的场景，而不可直接代表全域动机或人口。

将事件与地方层次连接时，应同时考察资源怎样被征集、信息怎样传递、迁徙怎样改变户籍和劳动、宗教与亲属网络怎样提供或限制协作。现有证据往往只照亮少数地点；保留这种不均衡，才能避免把宋辽夏金元的多政权历史改写成单一中心的必然过程。

账本条目\texttt{LIAO-0927-EVENT-029}所记927年的“耶律德光即位为辽太宗”，属于辽的历史语境。条目所示前因是前期边防、财政和统治联盟的压力构成直接背景，具体条件须按本条材料分辨，其可见后果为改变后续资源配置与政治选择；实际执行范围和社会影响仍须分项核验。这一变化若进入区域生活，必须通过道路、文书、军队、市场、家庭和地方中介才会发生；政权易手或法令发布不等于所有城镇、乡村与边地在同一时间获得同样经验。

本条所列来源为《辽史》〈本纪〉、〈营卫志〉、〈百官志〉；《宋史》辽国相关传；宋人使辽记录。它们能够支持特定的年序、制度语言、政治行动或局部遗存，却不能无条件化为社会总量。账本特别提醒“编年主干可靠；细节、规模与动机须与对方材料、地方文书或考古材料互校”。因此，军报的兵数、条约的名分、官府的族群分类、判牍中的道德判断以及后出的叙述都须回到产生它们的场景，而不可直接代表全域动机或人口。

将事件与地方层次连接时，应同时考察资源怎样被征集、信息怎样传递、迁徙怎样改变户籍和劳动、宗教与亲属网络怎样提供或限制协作。现有证据往往只照亮少数地点；保留这种不均衡，才能避免把宋辽夏金元的多政权历史改写成单一中心的必然过程。

\subsection{环境风险与地方经济}

账本条目\texttt{LIAO-0927-REGNAL-YEAR}所记927年的“辽以辽太宗在位第1年作为诏令、赋役与外交文书的纪年框架”，属于辽的历史语境。条目所示前因是新岁颁历、年号和纪年是中枢与地方行政沟通的共同前提，其可见后果为为同年政令、战事和地方材料提供日期锚点，不能单独推知具体政策执行。这一变化若进入区域生活，必须通过道路、文书、军队、市场、家庭和地方中介才会发生；政权易手或法令发布不等于所有城镇、乡村与边地在同一时间获得同样经验。

本条所列来源为《辽史》〈本纪〉、〈营卫志〉、〈百官志〉；《宋史》辽国相关传；宋人使辽记录。它们能够支持特定的年序、制度语言、政治行动或局部遗存，却不能无条件化为社会总量。账本特别提醒“纪年框架可靠；该记录仅确证政权纪年连续，年内具体事项须由本年条另证”。因此，军报的兵数、条约的名分、官府的族群分类、判牍中的道德判断以及后出的叙述都须回到产生它们的场景，而不可直接代表全域动机或人口。

将事件与地方层次连接时，应同时考察资源怎样被征集、信息怎样传递、迁徙怎样改变户籍和劳动、宗教与亲属网络怎样提供或限制协作。现有证据往往只照亮少数地点；保留这种不均衡，才能避免把宋辽夏金元的多政权历史改写成单一中心的必然过程。

账本条目\texttt{LIAO-0928-REGNAL-YEAR}所记928年的“辽以辽太宗在位第2年作为诏令、赋役与外交文书的纪年框架”，属于辽的历史语境。条目所示前因是新岁颁历、年号和纪年是中枢与地方行政沟通的共同前提，其可见后果为为同年政令、战事和地方材料提供日期锚点，不能单独推知具体政策执行。这一变化若进入区域生活，必须通过道路、文书、军队、市场、家庭和地方中介才会发生；政权易手或法令发布不等于所有城镇、乡村与边地在同一时间获得同样经验。

本条所列来源为《辽史》〈本纪〉、〈营卫志〉、〈百官志〉；《宋史》辽国相关传；宋人使辽记录。它们能够支持特定的年序、制度语言、政治行动或局部遗存，却不能无条件化为社会总量。账本特别提醒“纪年框架可靠；该记录仅确证政权纪年连续，年内具体事项须由本年条另证”。因此，军报的兵数、条约的名分、官府的族群分类、判牍中的道德判断以及后出的叙述都须回到产生它们的场景，而不可直接代表全域动机或人口。

将事件与地方层次连接时，应同时考察资源怎样被征集、信息怎样传递、迁徙怎样改变户籍和劳动、宗教与亲属网络怎样提供或限制协作。现有证据往往只照亮少数地点；保留这种不均衡，才能避免把宋辽夏金元的多政权历史改写成单一中心的必然过程。

账本条目\texttt{LIAO-0929-REGNAL-YEAR}所记929年的“辽以辽太宗在位第3年作为诏令、赋役与外交文书的纪年框架”，属于辽的历史语境。条目所示前因是新岁颁历、年号和纪年是中枢与地方行政沟通的共同前提，其可见后果为为同年政令、战事和地方材料提供日期锚点，不能单独推知具体政策执行。这一变化若进入区域生活，必须通过道路、文书、军队、市场、家庭和地方中介才会发生；政权易手或法令发布不等于所有城镇、乡村与边地在同一时间获得同样经验。

本条所列来源为《辽史》〈本纪〉、〈营卫志〉、〈百官志〉；《宋史》辽国相关传；宋人使辽记录。它们能够支持特定的年序、制度语言、政治行动或局部遗存，却不能无条件化为社会总量。账本特别提醒“纪年框架可靠；该记录仅确证政权纪年连续，年内具体事项须由本年条另证”。因此，军报的兵数、条约的名分、官府的族群分类、判牍中的道德判断以及后出的叙述都须回到产生它们的场景，而不可直接代表全域动机或人口。

将事件与地方层次连接时，应同时考察资源怎样被征集、信息怎样传递、迁徙怎样改变户籍和劳动、宗教与亲属网络怎样提供或限制协作。现有证据往往只照亮少数地点；保留这种不均衡，才能避免把宋辽夏金元的多政权历史改写成单一中心的必然过程。

账本条目\texttt{LIAO-0930-REGNAL-YEAR}所记930年的“辽以辽太宗在位第4年作为诏令、赋役与外交文书的纪年框架”，属于辽的历史语境。条目所示前因是新岁颁历、年号和纪年是中枢与地方行政沟通的共同前提，其可见后果为为同年政令、战事和地方材料提供日期锚点，不能单独推知具体政策执行。这一变化若进入区域生活，必须通过道路、文书、军队、市场、家庭和地方中介才会发生；政权易手或法令发布不等于所有城镇、乡村与边地在同一时间获得同样经验。

本条所列来源为《辽史》〈本纪〉、〈营卫志〉、〈百官志〉；《宋史》辽国相关传；宋人使辽记录。它们能够支持特定的年序、制度语言、政治行动或局部遗存，却不能无条件化为社会总量。账本特别提醒“纪年框架可靠；该记录仅确证政权纪年连续，年内具体事项须由本年条另证”。因此，军报的兵数、条约的名分、官府的族群分类、判牍中的道德判断以及后出的叙述都须回到产生它们的场景，而不可直接代表全域动机或人口。

将事件与地方层次连接时，应同时考察资源怎样被征集、信息怎样传递、迁徙怎样改变户籍和劳动、宗教与亲属网络怎样提供或限制协作。现有证据往往只照亮少数地点；保留这种不均衡，才能避免把宋辽夏金元的多政权历史改写成单一中心的必然过程。

账本条目\texttt{LIAO-0931-REGNAL-YEAR}所记931年的“辽以辽太宗在位第5年作为诏令、赋役与外交文书的纪年框架”，属于辽的历史语境。条目所示前因是新岁颁历、年号和纪年是中枢与地方行政沟通的共同前提，其可见后果为为同年政令、战事和地方材料提供日期锚点，不能单独推知具体政策执行。这一变化若进入区域生活，必须通过道路、文书、军队、市场、家庭和地方中介才会发生；政权易手或法令发布不等于所有城镇、乡村与边地在同一时间获得同样经验。

本条所列来源为《辽史》〈本纪〉、〈营卫志〉、〈百官志〉；《宋史》辽国相关传；宋人使辽记录。它们能够支持特定的年序、制度语言、政治行动或局部遗存，却不能无条件化为社会总量。账本特别提醒“纪年框架可靠；该记录仅确证政权纪年连续，年内具体事项须由本年条另证”。因此，军报的兵数、条约的名分、官府的族群分类、判牍中的道德判断以及后出的叙述都须回到产生它们的场景，而不可直接代表全域动机或人口。

将事件与地方层次连接时，应同时考察资源怎样被征集、信息怎样传递、迁徙怎样改变户籍和劳动、宗教与亲属网络怎样提供或限制协作。现有证据往往只照亮少数地点；保留这种不均衡，才能避免把宋辽夏金元的多政权历史改写成单一中心的必然过程。

账本条目\texttt{LIAO-0932-REGNAL-YEAR}所记932年的“辽以辽太宗在位第6年作为诏令、赋役与外交文书的纪年框架”，属于辽的历史语境。条目所示前因是新岁颁历、年号和纪年是中枢与地方行政沟通的共同前提，其可见后果为为同年政令、战事和地方材料提供日期锚点，不能单独推知具体政策执行。这一变化若进入区域生活，必须通过道路、文书、军队、市场、家庭和地方中介才会发生；政权易手或法令发布不等于所有城镇、乡村与边地在同一时间获得同样经验。

本条所列来源为《辽史》〈本纪〉、〈营卫志〉、〈百官志〉；《宋史》辽国相关传；宋人使辽记录。它们能够支持特定的年序、制度语言、政治行动或局部遗存，却不能无条件化为社会总量。账本特别提醒“纪年框架可靠；该记录仅确证政权纪年连续，年内具体事项须由本年条另证”。因此，军报的兵数、条约的名分、官府的族群分类、判牍中的道德判断以及后出的叙述都须回到产生它们的场景，而不可直接代表全域动机或人口。

将事件与地方层次连接时，应同时考察资源怎样被征集、信息怎样传递、迁徙怎样改变户籍和劳动、宗教与亲属网络怎样提供或限制协作。现有证据往往只照亮少数地点；保留这种不均衡，才能避免把宋辽夏金元的多政权历史改写成单一中心的必然过程。

账本条目\texttt{LIAO-0933-REGNAL-YEAR}所记933年的“辽以辽太宗在位第7年作为诏令、赋役与外交文书的纪年框架”，属于辽的历史语境。条目所示前因是新岁颁历、年号和纪年是中枢与地方行政沟通的共同前提，其可见后果为为同年政令、战事和地方材料提供日期锚点，不能单独推知具体政策执行。这一变化若进入区域生活，必须通过道路、文书、军队、市场、家庭和地方中介才会发生；政权易手或法令发布不等于所有城镇、乡村与边地在同一时间获得同样经验。

本条所列来源为《辽史》〈本纪〉、〈营卫志〉、〈百官志〉；《宋史》辽国相关传；宋人使辽记录。它们能够支持特定的年序、制度语言、政治行动或局部遗存，却不能无条件化为社会总量。账本特别提醒“纪年框架可靠；该记录仅确证政权纪年连续，年内具体事项须由本年条另证”。因此，军报的兵数、条约的名分、官府的族群分类、判牍中的道德判断以及后出的叙述都须回到产生它们的场景，而不可直接代表全域动机或人口。

将事件与地方层次连接时，应同时考察资源怎样被征集、信息怎样传递、迁徙怎样改变户籍和劳动、宗教与亲属网络怎样提供或限制协作。现有证据往往只照亮少数地点；保留这种不均衡，才能避免把宋辽夏金元的多政权历史改写成单一中心的必然过程。

账本条目\texttt{LIAO-0934-REGNAL-YEAR}所记934年的“辽以辽太宗在位第8年作为诏令、赋役与外交文书的纪年框架”，属于辽的历史语境。条目所示前因是新岁颁历、年号和纪年是中枢与地方行政沟通的共同前提，其可见后果为为同年政令、战事和地方材料提供日期锚点，不能单独推知具体政策执行。这一变化若进入区域生活，必须通过道路、文书、军队、市场、家庭和地方中介才会发生；政权易手或法令发布不等于所有城镇、乡村与边地在同一时间获得同样经验。

本条所列来源为《辽史》〈本纪〉、〈营卫志〉、〈百官志〉；《宋史》辽国相关传；宋人使辽记录。它们能够支持特定的年序、制度语言、政治行动或局部遗存，却不能无条件化为社会总量。账本特别提醒“纪年框架可靠；该记录仅确证政权纪年连续，年内具体事项须由本年条另证”。因此，军报的兵数、条约的名分、官府的族群分类、判牍中的道德判断以及后出的叙述都须回到产生它们的场景，而不可直接代表全域动机或人口。

将事件与地方层次连接时，应同时考察资源怎样被征集、信息怎样传递、迁徙怎样改变户籍和劳动、宗教与亲属网络怎样提供或限制协作。现有证据往往只照亮少数地点；保留这种不均衡，才能避免把宋辽夏金元的多政权历史改写成单一中心的必然过程。

账本条目\texttt{LIAO-0935-REGNAL-YEAR}所记935年的“辽以辽太宗在位第9年作为诏令、赋役与外交文书的纪年框架”，属于辽的历史语境。条目所示前因是新岁颁历、年号和纪年是中枢与地方行政沟通的共同前提，其可见后果为为同年政令、战事和地方材料提供日期锚点，不能单独推知具体政策执行。这一变化若进入区域生活，必须通过道路、文书、军队、市场、家庭和地方中介才会发生；政权易手或法令发布不等于所有城镇、乡村与边地在同一时间获得同样经验。

本条所列来源为《辽史》〈本纪〉、〈营卫志〉、〈百官志〉；《宋史》辽国相关传；宋人使辽记录。它们能够支持特定的年序、制度语言、政治行动或局部遗存，却不能无条件化为社会总量。账本特别提醒“纪年框架可靠；该记录仅确证政权纪年连续，年内具体事项须由本年条另证”。因此，军报的兵数、条约的名分、官府的族群分类、判牍中的道德判断以及后出的叙述都须回到产生它们的场景，而不可直接代表全域动机或人口。

将事件与地方层次连接时，应同时考察资源怎样被征集、信息怎样传递、迁徙怎样改变户籍和劳动、宗教与亲属网络怎样提供或限制协作。现有证据往往只照亮少数地点；保留这种不均衡，才能避免把宋辽夏金元的多政权历史改写成单一中心的必然过程。

账本条目\texttt{LIAO-0936-EVENT-030}所记936年的“辽支持石敬瑭建立后晋，十六州问题进入北方政治结构”，属于辽与后晋的历史语境。条目所示前因是前期边防、财政和统治联盟的压力构成直接背景，具体条件须按本条材料分辨，其可见后果为改变后续资源配置与政治选择；实际执行范围和社会影响仍须分项核验。这一变化若进入区域生活，必须通过道路、文书、军队、市场、家庭和地方中介才会发生；政权易手或法令发布不等于所有城镇、乡村与边地在同一时间获得同样经验。

本条所列来源为《辽史》〈本纪〉、〈营卫志〉、〈百官志〉；《宋史》辽国相关传；宋人使辽记录。它们能够支持特定的年序、制度语言、政治行动或局部遗存，却不能无条件化为社会总量。账本特别提醒“编年主干可靠；细节、规模与动机须与对方材料、地方文书或考古材料互校”。因此，军报的兵数、条约的名分、官府的族群分类、判牍中的道德判断以及后出的叙述都须回到产生它们的场景，而不可直接代表全域动机或人口。

将事件与地方层次连接时，应同时考察资源怎样被征集、信息怎样传递、迁徙怎样改变户籍和劳动、宗教与亲属网络怎样提供或限制协作。现有证据往往只照亮少数地点；保留这种不均衡，才能避免把宋辽夏金元的多政权历史改写成单一中心的必然过程。

账本条目\texttt{LIAO-0936-REGNAL-YEAR}所记936年的“辽以辽太宗在位第10年作为诏令、赋役与外交文书的纪年框架”，属于辽的历史语境。条目所示前因是新岁颁历、年号和纪年是中枢与地方行政沟通的共同前提，其可见后果为为同年政令、战事和地方材料提供日期锚点，不能单独推知具体政策执行。这一变化若进入区域生活，必须通过道路、文书、军队、市场、家庭和地方中介才会发生；政权易手或法令发布不等于所有城镇、乡村与边地在同一时间获得同样经验。

本条所列来源为《辽史》〈本纪〉、〈营卫志〉、〈百官志〉；《宋史》辽国相关传；宋人使辽记录。它们能够支持特定的年序、制度语言、政治行动或局部遗存，却不能无条件化为社会总量。账本特别提醒“纪年框架可靠；该记录仅确证政权纪年连续，年内具体事项须由本年条另证”。因此，军报的兵数、条约的名分、官府的族群分类、判牍中的道德判断以及后出的叙述都须回到产生它们的场景，而不可直接代表全域动机或人口。

将事件与地方层次连接时，应同时考察资源怎样被征集、信息怎样传递、迁徙怎样改变户籍和劳动、宗教与亲属网络怎样提供或限制协作。现有证据往往只照亮少数地点；保留这种不均衡，才能避免把宋辽夏金元的多政权历史改写成单一中心的必然过程。

账本条目\texttt{LIAO-0937-REGNAL-YEAR}所记937年的“辽以辽太宗在位第11年作为诏令、赋役与外交文书的纪年框架”，属于辽的历史语境。条目所示前因是新岁颁历、年号和纪年是中枢与地方行政沟通的共同前提，其可见后果为为同年政令、战事和地方材料提供日期锚点，不能单独推知具体政策执行。这一变化若进入区域生活，必须通过道路、文书、军队、市场、家庭和地方中介才会发生；政权易手或法令发布不等于所有城镇、乡村与边地在同一时间获得同样经验。

本条所列来源为《辽史》〈本纪〉、〈营卫志〉、〈百官志〉；《宋史》辽国相关传；宋人使辽记录。它们能够支持特定的年序、制度语言、政治行动或局部遗存，却不能无条件化为社会总量。账本特别提醒“纪年框架可靠；该记录仅确证政权纪年连续，年内具体事项须由本年条另证”。因此，军报的兵数、条约的名分、官府的族群分类、判牍中的道德判断以及后出的叙述都须回到产生它们的场景，而不可直接代表全域动机或人口。

将事件与地方层次连接时，应同时考察资源怎样被征集、信息怎样传递、迁徙怎样改变户籍和劳动、宗教与亲属网络怎样提供或限制协作。现有证据往往只照亮少数地点；保留这种不均衡，才能避免把宋辽夏金元的多政权历史改写成单一中心的必然过程。

\subsection{宗教、法律与材料边界}

账本条目\texttt{LIAO-0938-EVENT-031}所记938年的“辽以幽州为南京，经营燕云城市与交通节点”，属于辽的历史语境。条目所示前因是前期边防、财政和统治联盟的压力构成直接背景，具体条件须按本条材料分辨，其可见后果为改变后续资源配置与政治选择；实际执行范围和社会影响仍须分项核验。这一变化若进入区域生活，必须通过道路、文书、军队、市场、家庭和地方中介才会发生；政权易手或法令发布不等于所有城镇、乡村与边地在同一时间获得同样经验。

本条所列来源为《辽史》〈本纪〉、〈营卫志〉、〈百官志〉；《宋史》辽国相关传；宋人使辽记录。它们能够支持特定的年序、制度语言、政治行动或局部遗存，却不能无条件化为社会总量。账本特别提醒“编年主干可靠；细节、规模与动机须与对方材料、地方文书或考古材料互校”。因此，军报的兵数、条约的名分、官府的族群分类、判牍中的道德判断以及后出的叙述都须回到产生它们的场景，而不可直接代表全域动机或人口。

将事件与地方层次连接时，应同时考察资源怎样被征集、信息怎样传递、迁徙怎样改变户籍和劳动、宗教与亲属网络怎样提供或限制协作。现有证据往往只照亮少数地点；保留这种不均衡，才能避免把宋辽夏金元的多政权历史改写成单一中心的必然过程。

账本条目\texttt{LIAO-0938-REGNAL-YEAR}所记938年的“辽以辽太宗在位第12年作为诏令、赋役与外交文书的纪年框架”，属于辽的历史语境。条目所示前因是新岁颁历、年号和纪年是中枢与地方行政沟通的共同前提，其可见后果为为同年政令、战事和地方材料提供日期锚点，不能单独推知具体政策执行。这一变化若进入区域生活，必须通过道路、文书、军队、市场、家庭和地方中介才会发生；政权易手或法令发布不等于所有城镇、乡村与边地在同一时间获得同样经验。

本条所列来源为《辽史》〈本纪〉、〈营卫志〉、〈百官志〉；《宋史》辽国相关传；宋人使辽记录。它们能够支持特定的年序、制度语言、政治行动或局部遗存，却不能无条件化为社会总量。账本特别提醒“纪年框架可靠；该记录仅确证政权纪年连续，年内具体事项须由本年条另证”。因此，军报的兵数、条约的名分、官府的族群分类、判牍中的道德判断以及后出的叙述都须回到产生它们的场景，而不可直接代表全域动机或人口。

将事件与地方层次连接时，应同时考察资源怎样被征集、信息怎样传递、迁徙怎样改变户籍和劳动、宗教与亲属网络怎样提供或限制协作。现有证据往往只照亮少数地点；保留这种不均衡，才能避免把宋辽夏金元的多政权历史改写成单一中心的必然过程。

账本条目\texttt{LIAO-0939-REGNAL-YEAR}所记939年的“辽以辽太宗在位第13年作为诏令、赋役与外交文书的纪年框架”，属于辽的历史语境。条目所示前因是新岁颁历、年号和纪年是中枢与地方行政沟通的共同前提，其可见后果为为同年政令、战事和地方材料提供日期锚点，不能单独推知具体政策执行。这一变化若进入区域生活，必须通过道路、文书、军队、市场、家庭和地方中介才会发生；政权易手或法令发布不等于所有城镇、乡村与边地在同一时间获得同样经验。

本条所列来源为《辽史》〈本纪〉、〈营卫志〉、〈百官志〉；《宋史》辽国相关传；宋人使辽记录。它们能够支持特定的年序、制度语言、政治行动或局部遗存，却不能无条件化为社会总量。账本特别提醒“纪年框架可靠；该记录仅确证政权纪年连续，年内具体事项须由本年条另证”。因此，军报的兵数、条约的名分、官府的族群分类、判牍中的道德判断以及后出的叙述都须回到产生它们的场景，而不可直接代表全域动机或人口。

将事件与地方层次连接时，应同时考察资源怎样被征集、信息怎样传递、迁徙怎样改变户籍和劳动、宗教与亲属网络怎样提供或限制协作。现有证据往往只照亮少数地点；保留这种不均衡，才能避免把宋辽夏金元的多政权历史改写成单一中心的必然过程。

账本条目\texttt{LIAO-0940-REGNAL-YEAR}所记940年的“辽以辽太宗在位第14年作为诏令、赋役与外交文书的纪年框架”，属于辽的历史语境。条目所示前因是新岁颁历、年号和纪年是中枢与地方行政沟通的共同前提，其可见后果为为同年政令、战事和地方材料提供日期锚点，不能单独推知具体政策执行。这一变化若进入区域生活，必须通过道路、文书、军队、市场、家庭和地方中介才会发生；政权易手或法令发布不等于所有城镇、乡村与边地在同一时间获得同样经验。

本条所列来源为《辽史》〈本纪〉、〈营卫志〉、〈百官志〉；《宋史》辽国相关传；宋人使辽记录。它们能够支持特定的年序、制度语言、政治行动或局部遗存，却不能无条件化为社会总量。账本特别提醒“纪年框架可靠；该记录仅确证政权纪年连续，年内具体事项须由本年条另证”。因此，军报的兵数、条约的名分、官府的族群分类、判牍中的道德判断以及后出的叙述都须回到产生它们的场景，而不可直接代表全域动机或人口。

将事件与地方层次连接时，应同时考察资源怎样被征集、信息怎样传递、迁徙怎样改变户籍和劳动、宗教与亲属网络怎样提供或限制协作。现有证据往往只照亮少数地点；保留这种不均衡，才能避免把宋辽夏金元的多政权历史改写成单一中心的必然过程。

账本条目\texttt{LIAO-0941-REGNAL-YEAR}所记941年的“辽以辽太宗在位第15年作为诏令、赋役与外交文书的纪年框架”，属于辽的历史语境。条目所示前因是新岁颁历、年号和纪年是中枢与地方行政沟通的共同前提，其可见后果为为同年政令、战事和地方材料提供日期锚点，不能单独推知具体政策执行。这一变化若进入区域生活，必须通过道路、文书、军队、市场、家庭和地方中介才会发生；政权易手或法令发布不等于所有城镇、乡村与边地在同一时间获得同样经验。

本条所列来源为《辽史》〈本纪〉、〈营卫志〉、〈百官志〉；《宋史》辽国相关传；宋人使辽记录。它们能够支持特定的年序、制度语言、政治行动或局部遗存，却不能无条件化为社会总量。账本特别提醒“纪年框架可靠；该记录仅确证政权纪年连续，年内具体事项须由本年条另证”。因此，军报的兵数、条约的名分、官府的族群分类、判牍中的道德判断以及后出的叙述都须回到产生它们的场景，而不可直接代表全域动机或人口。

将事件与地方层次连接时，应同时考察资源怎样被征集、信息怎样传递、迁徙怎样改变户籍和劳动、宗教与亲属网络怎样提供或限制协作。现有证据往往只照亮少数地点；保留这种不均衡，才能避免把宋辽夏金元的多政权历史改写成单一中心的必然过程。

账本条目\texttt{LIAO-0942-REGNAL-YEAR}所记942年的“辽以辽太宗在位第16年作为诏令、赋役与外交文书的纪年框架”，属于辽的历史语境。条目所示前因是新岁颁历、年号和纪年是中枢与地方行政沟通的共同前提，其可见后果为为同年政令、战事和地方材料提供日期锚点，不能单独推知具体政策执行。这一变化若进入区域生活，必须通过道路、文书、军队、市场、家庭和地方中介才会发生；政权易手或法令发布不等于所有城镇、乡村与边地在同一时间获得同样经验。

本条所列来源为《辽史》〈本纪〉、〈营卫志〉、〈百官志〉；《宋史》辽国相关传；宋人使辽记录。它们能够支持特定的年序、制度语言、政治行动或局部遗存，却不能无条件化为社会总量。账本特别提醒“纪年框架可靠；该记录仅确证政权纪年连续，年内具体事项须由本年条另证”。因此，军报的兵数、条约的名分、官府的族群分类、判牍中的道德判断以及后出的叙述都须回到产生它们的场景，而不可直接代表全域动机或人口。

将事件与地方层次连接时，应同时考察资源怎样被征集、信息怎样传递、迁徙怎样改变户籍和劳动、宗教与亲属网络怎样提供或限制协作。现有证据往往只照亮少数地点；保留这种不均衡，才能避免把宋辽夏金元的多政权历史改写成单一中心的必然过程。

账本条目\texttt{LIAO-0943-REGNAL-YEAR}所记943年的“辽以辽太宗在位第17年作为诏令、赋役与外交文书的纪年框架”，属于辽的历史语境。条目所示前因是新岁颁历、年号和纪年是中枢与地方行政沟通的共同前提，其可见后果为为同年政令、战事和地方材料提供日期锚点，不能单独推知具体政策执行。这一变化若进入区域生活，必须通过道路、文书、军队、市场、家庭和地方中介才会发生；政权易手或法令发布不等于所有城镇、乡村与边地在同一时间获得同样经验。

本条所列来源为《辽史》〈本纪〉、〈营卫志〉、〈百官志〉；《宋史》辽国相关传；宋人使辽记录。它们能够支持特定的年序、制度语言、政治行动或局部遗存，却不能无条件化为社会总量。账本特别提醒“纪年框架可靠；该记录仅确证政权纪年连续，年内具体事项须由本年条另证”。因此，军报的兵数、条约的名分、官府的族群分类、判牍中的道德判断以及后出的叙述都须回到产生它们的场景，而不可直接代表全域动机或人口。

将事件与地方层次连接时，应同时考察资源怎样被征集、信息怎样传递、迁徙怎样改变户籍和劳动、宗教与亲属网络怎样提供或限制协作。现有证据往往只照亮少数地点；保留这种不均衡，才能避免把宋辽夏金元的多政权历史改写成单一中心的必然过程。

账本条目\texttt{LIAO-0944-REGNAL-YEAR}所记944年的“辽以辽太宗在位第18年作为诏令、赋役与外交文书的纪年框架”，属于辽的历史语境。条目所示前因是新岁颁历、年号和纪年是中枢与地方行政沟通的共同前提，其可见后果为为同年政令、战事和地方材料提供日期锚点，不能单独推知具体政策执行。这一变化若进入区域生活，必须通过道路、文书、军队、市场、家庭和地方中介才会发生；政权易手或法令发布不等于所有城镇、乡村与边地在同一时间获得同样经验。

本条所列来源为《辽史》〈本纪〉、〈营卫志〉、〈百官志〉；《宋史》辽国相关传；宋人使辽记录。它们能够支持特定的年序、制度语言、政治行动或局部遗存，却不能无条件化为社会总量。账本特别提醒“纪年框架可靠；该记录仅确证政权纪年连续，年内具体事项须由本年条另证”。因此，军报的兵数、条约的名分、官府的族群分类、判牍中的道德判断以及后出的叙述都须回到产生它们的场景，而不可直接代表全域动机或人口。

将事件与地方层次连接时，应同时考察资源怎样被征集、信息怎样传递、迁徙怎样改变户籍和劳动、宗教与亲属网络怎样提供或限制协作。现有证据往往只照亮少数地点；保留这种不均衡，才能避免把宋辽夏金元的多政权历史改写成单一中心的必然过程。

账本条目\texttt{LIAO-0945-REGNAL-YEAR}所记945年的“辽以辽太宗在位第19年作为诏令、赋役与外交文书的纪年框架”，属于辽的历史语境。条目所示前因是新岁颁历、年号和纪年是中枢与地方行政沟通的共同前提，其可见后果为为同年政令、战事和地方材料提供日期锚点，不能单独推知具体政策执行。这一变化若进入区域生活，必须通过道路、文书、军队、市场、家庭和地方中介才会发生；政权易手或法令发布不等于所有城镇、乡村与边地在同一时间获得同样经验。

本条所列来源为《辽史》〈本纪〉、〈营卫志〉、〈百官志〉；《宋史》辽国相关传；宋人使辽记录。它们能够支持特定的年序、制度语言、政治行动或局部遗存，却不能无条件化为社会总量。账本特别提醒“纪年框架可靠；该记录仅确证政权纪年连续，年内具体事项须由本年条另证”。因此，军报的兵数、条约的名分、官府的族群分类、判牍中的道德判断以及后出的叙述都须回到产生它们的场景，而不可直接代表全域动机或人口。

将事件与地方层次连接时，应同时考察资源怎样被征集、信息怎样传递、迁徙怎样改变户籍和劳动、宗教与亲属网络怎样提供或限制协作。现有证据往往只照亮少数地点；保留这种不均衡，才能避免把宋辽夏金元的多政权历史改写成单一中心的必然过程。
